\documentclass[11pt, letterpaper]{article}
\usepackage{amsmath, amssymb}
\usepackage{hyperref}
\usepackage{booktabs}
\usepackage{enumitem}
\usepackage{xcolor}
\usepackage{geometry}
\usepackage{fancyhdr}
\usepackage{titlesec}
\usepackage{tcolorbox}
\usepackage{parskip}
\geometry{margin=1in}

\definecolor{courseblue}{HTML}{002452}
\definecolor{coursegold}{HTML}{FABD0F}
\definecolor{coursered}{HTML}{B90E31}
\definecolor{lightblue}{HTML}{E8EEF7}

\hypersetup{colorlinks=true,linkcolor=courseblue,urlcolor=coursered}

\tcbuselibrary{skins}
\newtcolorbox{guidancebox}{enhanced,colback=lightblue,colframe=courseblue,
  sharp corners,leftrule=4pt,left=6pt,right=6pt,top=4pt,bottom=4pt}

\pagestyle{fancy}\fancyhf{}
\fancyhead[L]{\small\textcolor{courseblue}{\textbf{CISC 473 --- Deep Learning | Fall 2026}}}
\fancyhead[R]{\small\textcolor{courseblue}{Project Proposal}}
\fancyfoot[C]{\small\thepage}
\renewcommand{\headrulewidth}{0.4pt}

\titleformat{\section}{\large\bfseries\color{courseblue}}{}{0em}{}[\titlerule]
\titleformat{\subsection}{\normalsize\bfseries\color{courseblue}}{}{0em}{}

%% ════════════════════════════════════════════════════════════════
%% INSTRUCTIONS FOR STUDENTS
%% Replace every [PLACEHOLDER] with your actual content.
%% Delete all \guidancebox blocks before submitting.
%% Submit as a compiled PDF + .tex source in a .zip archive.
%% ════════════════════════════════════════════════════════════════

\begin{document}

{\centering
  {\Large\textbf{CISC 473: Deep Learning --- Project Proposal}}\\[0.5em]
  {\large\textit{P22: Cross-Architecture Knowledge Distillation}}\\[0.4em]
  \textbf{Group Members:}\\{}
  Lee Xuan \quad 20630324 \quad (Role: Engineering Lead)\\{}
  Fong Yee Xin \quad [Student ID] \quad (Role: Engineering Lead)\\[0.3em]
  \textbf{GitHub Repository:} \url{https://github.com/l33zard/cakd.git}\\[0.3em]
  \textbf{Date:} \today\\[0.4em]
  \rule{\linewidth}{0.4pt}
}

\section{Problem Statement and Motivation}

\begin{guidancebox}
\textbf{Guidance (delete before submitting).}
Clearly define the problem you are solving. Explain \emph{why} it matters
in a broader context. Reference 1--2 high-level papers or application areas.
This section should answer: ``What problem are we addressing, and who cares?''
Target: $\approx$300 words.
\end{guidancebox}

[WRITE YOUR PROBLEM STATEMENT HERE. Describe the task, why it is hard, and why
it is important. Situate it within a research or application context.]

Knowledge distillation is the process of compressing a large and more complex model, known as the teacher, into a smaller, simpler model, known as the student, and in the process of doing so, preserve much of the teacher's accuracy. This process is relevant in the current day and age for 3 different reasons. 

Firstly, running a model for everyday users translates to companies needing to route billions of daily user queries through massive, energy-hungry data centers, an expensive process. Through distillation, the model can be shrunk such that they require significantly less computing power, minimising cloud hosting and electricity bills. 

Secondly, there is currently a transition from cloud-based AI to AI systems that run directly on personal hardware. This would allow AI features to work instantly, offline and without uploading sensitive data to an external server. This is however limited by the strict memory and battery constraints that exist in smaller personal devices like smartphones, smartwatches, laptops and automotive computer chips. Knowledge distillation would compress state-of-the-art AI so that these personal devices could better handle AI systems directly. 

Thirdly, knowledge distillation enables smaller student models to retain much of the predictive performance of larger teacher models while requiring less computation during inference. This can reduce inference latency and resource requirements, making distilled models more suitable for real-time and resource-constrained applications such as autonomous systems, robotics, and edge devices, fields that are increasingly adopting AI-driven systems. 

However, knowledge distillation becomes more challenging when the teacher and student use different architectures. In particular, we look at the architecture gap between convolutional neural network (CNN) and vision transformer (ViT) as they learn and represent visual information differently, making the transfer of intermediate representations between them more complex. To be more specific, CNNs represent images through hierarchical spatial feature maps produced by convolutional operations, while ViTs represent images as sequences of patch tokens and model relationships between them using self-attention. As a result, their intermediate representations differ in structure, making direct feature-level knowledge transfer difficult. This project investigates different methods of knowledge distillation, benchmarked on CIFAR-100, and proposes an intermediate-feature alignment strategy for CNN-to-ViT knowledge distillation. 
\section{Literature Survey}

\begin{guidancebox}
\textbf{Guidance (delete before submitting).}
Summarise at least \textbf{5 related papers}. For each paper:
(1) state the core idea in one sentence, (2) describe what it achieves,
and (3) explain how your project relates to or differs from it.
Avoid copying abstracts verbatim. Identify at least one open problem
that motivates your proposed approach.
\end{guidancebox}

\subsection{Related Works}

\paragraph{[Decoupled Knowledge Distillation: Borui Zhao, Quan Cui, Renjie Song, Yiyu Qiu, Jiajun Liang, 2022, CVPR 2022.]}
[One-paragraph summary: what they proposed, what benchmark they evaluated on,
what metric they achieved, and how this relates to your project.]

\paragraph{[Contrastive Representation Distillation: Yonglong Tian, Dilip Krishnan, Phillip Isola, 2020, ICLR 2020.]}
[Summary\ldots]

\paragraph{[Distilling Knowledge via Knowledge Review: Pengguang Chen, Shu Liu, Hengshuang Zhao, Jiaya Jia, 2021, CVPR 2021.]}
[Summary\ldots]

\paragraph{[Paper 4: Authors, Year, Venue.]}
[Summary\ldots]

\paragraph{[Paper 5: Authors, Year, Venue.]}
[Summary\ldots]

\subsection{Open Problem}

\begin{guidancebox}
\textbf{Guidance.} Identify a specific gap, limitation, or unexplored direction
in the related work that motivates your novel contribution. Be precise.
\end{guidancebox}

[DESCRIBE THE OPEN PROBLEM OR GAP YOUR PROJECT WILL ADDRESS.]

\section{Proposed Methodology}

\begin{guidancebox}
\textbf{Guidance.}
Describe your approach in enough detail that it could be reproduced:
(1) baseline model and dataset you will start with,
(2) your proposed modification or extension and \emph{why} you chose it,
(3) how you will evaluate success (metrics, baselines).
Include a diagram if helpful. Target: $\approx$400 words.
\end{guidancebox}

\subsection{Baseline}

[DESCRIBE THE BASELINE MODEL AND DATASET. Which existing code/checkpoint will
you start from? Which metric will you use to evaluate the baseline?]

We will use CIFAR-100 as the primary dataset for training and evaluating all models. We will first establish non-distilled performance by evaluating the teacher and student models using standard supervised training. For the main cross-architecture experiment, a ResNet-50 will be used as the CNN teacher and a DeiT-Tiny as the ViT student. The performance of the DeiT-Tiny trained without knowledge distillation will serve as the primary student baseline.
We will then reproduce established knowledge distillation methods, including Vanilla Knowledge Distillation (KD), Contrastive Representation Distillation (CRD), and Decoupled Knowledge Distillation (DKD). These methods will provide additional baselines for determining whether our proposed cross-architecture feature-alignment method provides an improvement over existing distillation approaches. We will build upon existing official or publicly available implementations and pretrained checkpoints where appropriate. The primary evaluation metric will be top-1 classification accuracy on the CIFAR-100 test set, with computational cost measured using FLOPs.

\subsection{Proposed Extension}

[DESCRIBE YOUR PROPOSED MODIFICATION OR EXTENSION. Why do you expect it to
help? What is the key hypothesis you are testing?]

Our proposed extension focuses on improving the transfer of intermediate representations between a CNN teacher and a ViT student. Because these representations differ in both structure and dimensionality, they cannot be directly aligned for feature-level knowledge distillation.
We propose an intermediate-feature alignment module that transforms the CNN teacher's feature maps into a representation compatible with the ViT student's patch-token representations. The CNN feature maps will first be converted from their spatial representation into a sequence of feature vectors. A learnable projection will then transform the CNN channel dimension to the ViT embedding dimension, while spatial pooling or interpolation can be used where necessary to align the number of teacher features with the number of student patch tokens. This enables a feature-level distillation loss to be computed between the aligned teacher and student representations.
Our key hypothesis is that explicitly aligning the intermediate representations of the CNN teacher and ViT student will allow the student to learn additional information that is not captured solely by the teacher's output predictions, thereby improving cross-architecture knowledge distillation.

\subsection{Evaluation Plan}

[DESCRIBE YOUR EVALUATION PLAN:
\begin{itemize}[noitemsep]
  \item Dataset(s) and split(s) you will use
  \item Metrics (e.g., PSNR/SSIM, ROUGE-L, F1, AUROC)
  \item Baselines you will compare against
  \item Number of seeds / runs for statistical reliability
\end{itemize}]

All experiments will use CIFAR-100 with consistent training, validation, and test splits across methods. We will compare the normally trained DeiT-Tiny student against Vanilla KD, CRD, DKD, and our proposed intermediate-feature alignment method using a ResNet-50 teacher.
The primary metric will be top-1 classification accuracy, while FLOPs will be reported to evaluate computational efficiency. We will additionally use Centered Kernel Alignment (CKA) to measure the similarity between teacher and student intermediate representations before and after distillation. This will allow us to investigate whether our proposed method produces greater representational alignment between the CNN and ViT.
We will perform ablation studies over the distillation temperature \(T \in \{1,2,4,8,16\}\) and the weights assigned to the different components of the distillation loss. Each main experiment will be conducted using at least three independent runs with fixed random seeds, and the resulting performance will be aggregated to improve statistical reliability.

\subsection{Early Novel Contribution (Optional Bonus)}

\begin{guidancebox}
\textbf{Guidance.} If you have already identified a specific novel contribution
(a gap in existing evaluations, an untested configuration, a new domain),
describe it here for the +2 bonus points. Be concrete and specific.
\end{guidancebox}

[DESCRIBE YOUR NOVEL CONTRIBUTION IDEA, OR LEAVE THIS SECTION EMPTY.]

\section{Project Timeline and Roles}

\begin{guidancebox}
\textbf{Guidance.} Provide a week-by-week plan. Assign each task to a named
group member. The plan must be specific enough to hold yourselves accountable.
\end{guidancebox}

\begin{center}
\begin{tabular}{lp{7cm}l}
\toprule
\textbf{Week} & \textbf{Tasks} & \textbf{Responsible} \\
\midrule
1--2 & Literature review; CIFAR-100 pipeline; teacher and student baseline training (ResNet-110/MobileNetV2); \textbf{Proposal due} & Fong Yee Xin \\
3 & Vanilla KD; temperature and distillation-weight ablations & Lee Xuan \\
4 & CRD and DKD distillation baselines & Fong Yee Xin \\
5 & ResNet-50 and DeiT-Tiny baselines; cross-architecture distillation: ResNet-50 $\rightarrow$ DeiT-Tiny & Lee Xuan \\
6 & CKA analysis before/after distillation; \textbf{Midterm due} & Fong Yee Xin \\
7 & Develop intermediate-feature alignment for CNN $\rightarrow$ ViT distillation & Lee Xuan \\
8 & Extended experiments and ablations; efficiency analysis (accuracy vs.\ FLOPs) & Fong Yee Xin \\
9 & Results analysis; accuracy tables; CKA heatmaps; report writing & Lee Xuan \\
10 & Final report and reproducibility checks; \textbf{Final due} & Fong Yee Xin \\
\bottomrule
\end{tabular}
\end{center}

\subsection{Contribution Declaration}

\begin{tabular}{lcc}
\toprule
\textbf{Member} & \textbf{Primary Responsibilities} & \textbf{Estimated \%} \\
\midrule
{Lee Xuan} & Literature review, report writing, Data pipeline, evaluation code & 50\%  \\
{Fong Yee Xin}& README, Model implementation, experiments & 50\%  \\
\bottomrule
\end{tabular}

\section{References}

\begin{guidancebox}
\textbf{Guidance.} Use consistent citation formatting. Include at minimum the 5
papers from Section~2, plus any other papers you reference.
arXiv IDs and DOIs are preferred over URLs.
\end{guidancebox}

\begin{enumerate}[label={[\arabic*]},leftmargin=*]
  \item [Author(s), Year. ``Title.'' \textit{Venue}. arXiv:XXXX.XXXXX or DOI.]
  \item [Author(s), Year. ``Title.'' \textit{Venue}.]
  \item [\ldots]
\end{enumerate}

\end{document}
